% Generated draft: requires mathematical and source review.
\documentclass[11pt]{article}
\usepackage[margin=25mm]{geometry}
\usepackage{fontspec}
\setmainfont{Latin Modern Roman}
\usepackage{amsmath,amssymb,mathtools,unicode-math}
\setmathfont{Latin Modern Math}
\usepackage{enumitem,needspace,xurl,hyperref}
\setlength{\parindent}{0pt}
\setlength{\parskip}{5pt}
\setlength{\emergencystretch}{3em}
\newcommand{\N}{\mathbb{N}}
\newcommand{\Z}{\mathbb{Z}}
\newcommand{\Q}{\mathbb{Q}}
\newcommand{\R}{\mathbb{R}}
\newcommand{\C}{\mathbb{C}}
\newcommand{\F}{\mathbb{F}}
\newcommand{\A}{\mathbb{A}}
\newcommand{\PP}{\mathbb{P}}
\newcommand{\E}{\mathbb{E}}
\newcommand{\eps}{\varepsilon}
\newcommand{\dd}{\,\mathrm{d}}
\begin{document}
\section*{UnsolvedMath ID: 30001878}
\textbf{Source problem code:} catalogue entry
\subsection{Definitions and mathematical statement}
For positive integers \(a_1,\ldots,a_k\), with \(k\geq 1\), define the finite regular continued fraction by
\[
[0;a_1,\ldots,a_k]
=\cfrac{1}{a_1+\cfrac{1}{a_2+\cfrac{1}{\ddots+\cfrac{1}{a_k}}}},
\qquad [0;a_1]=\frac{1}{a_1}.
\]
The denominator of a positive rational number means its positive denominator when written in lowest terms.

\textbf{Question (Zaremba's bounded continued-fraction conjecture, as identified by the supplied record).} Does there exist an integer \(A\geq 1\) such that, for every integer \(q\geq 1\), there exist integers \(k\geq 1\), \(p\) with \(1\leq p\leq q\), and \(a_1,\ldots,a_k\) satisfying
\[
\gcd(p,q)=1,\qquad
1\leq a_i\leq A\quad(1\leq i\leq k),\qquad
\frac{p}{q}=[0;a_1,\ldots,a_k]?
\]
Here \(A\) is independent of \(q\), while the numerator, length, and partial quotients may depend on \(q\). The boundary case \(q=1\) is included: the expansion \([0;1]=1\) is allowed. No condition \(a_k\geq 2\) is imposed.

The source writes \([a_1,\ldots,a_k]\) without specifying whether the initial reciprocal is implicit. The display above uses the reciprocal convention explicitly. If the source instead means the integer-part convention \([a_1;a_2,\ldots,a_k]\), its denominator for \(k\geq 2\) equals the reduced denominator of \([0;a_2,\ldots,a_k]\); conversely, prepending the integer part \(1\) preserves the denominator of any reciprocal expansion. Thus the existential denominator-coverage question is equivalent under these two conventions when lengths are unrestricted. The intended source notation still requires confirmation.
\subsection{Short English statement}
Is there a single positive integer bound \(A\) such that every positive integer occurs as the reduced denominator of a finite regular continued fraction whose positive partial quotients are all at most \(A\)? The supplied record identifies this as Zaremba's conjecture and labels it partially solved; that label is supplied dated curation, not an independently verified mathematical assessment.
\subsection{Sources}
\begin{itemize}
\item Supplied local UnsolvedMath registry record: permanent ID \(30001878\), exact source problem code \texttt{catalogue entry}, title ``Zaremba's Bounded Continued-Fraction Conjecture,'' in the Oberwolfach Reports Open Problems collection. The record gives proposed year \(2011\) and supplies no named proposer.
\item Supplied report citation: \emph{The Analytic Theory of Automorphic Forms} (2012), Oberwolfach Reports, pp.\ 2397--2458, DOI \texttt{10.4171/owr/2011/43}. Source URL and literature-source label ``original OWR report'': \url{https://doi.org/10.4171/owr/2011/43}. The supplied background describes these reports as documenting mathematical workshops and their open-problem discussions.
\item Supplied literature source: ``Zaremba's conjecture for geometric sequences,'' \url{https://www.tandfonline.com/doi/abs/10.1080/00029890.2024.2323902}.
\item Supplied literature source: ``radical bound for Zaremba's conjecture,'' \url{https://doi.org/10.1112/blms.13087}.
\item Literature assessment dated exactly \texttt{2026-08-21}: the supplied curation labels the problem partially solved, says full bounded-partial-quotient coverage remains open, and attributes substantial density and special-family results to Bourgain--Kontorovich and later work. This is supplied dated curation derived from the OWR material, not independently verified here; the record itself requests independent verification of time-sensitive status and cited claims.
\item The supplied statement status is \texttt{corrected\_verified}. Its verification note reports that the mathematical content was checked against the official Oberwolfach report and that notation, context, and question form were normalized without changing the source-level claim. This records the supplied verification claim; no source visit or independent verification is asserted here.
\end{itemize}
\end{document}
